\documentclass[10pt,a4paper]{amsart}
\usepackage{amsmath,amssymb,amsthm,graphicx}
% Retained preamble for the Stacks Project chapter {CH} (shared project preamble).
% The chapter file pinned-original-source.tex only carries \input{{preamble}} on its first line,
% so no preamble text lives inside the chapter itself.  The Stacks Project uses one single shared
% preamble for every chapter: the snapshot stages it three times (sources/stacks-chow/preamble.tex,
% sources/stacks-etale-cohomology/preamble.tex and focused-seeds/AGP/sources/stacks-snapshot-20261011/
% preamble.tex) and those copies are identical in size.  The block below is copied mechanically,
% line for line, from sources/stacks-chow/preamble.tex lines 14-18 and 154-217: the Xy-pic setup and
% the native theorem declarations together with the mathematical macros.  Nothing was added, removed
% or redefined; the wrapper class and the amsmath/amssymb/amsthm/graphicx base packages are supplied
% by the assembler.
% For commutative diagrams we use Xy-pic
% Chapter: stacks-spaces-more-morphisms.tex -- More on Morphisms of Spaces.
\usepackage[all]{xy}

% We use 2cell for 2-commutative diagrams.
\xyoption{2cell}

% Theorem environments.
%
\theoremstyle{plain}
\newtheorem{theorem}[subsection]{Theorem}
\newtheorem{proposition}[subsection]{Proposition}
\newtheorem{lemma}[subsection]{Lemma}

\theoremstyle{definition}
\newtheorem{definition}[subsection]{Definition}
\newtheorem{example}[subsection]{Example}
\newtheorem{exercise}[subsection]{Exercise}
\newtheorem{situation}[subsection]{Situation}

\theoremstyle{remark}
\newtheorem{remark}[subsection]{Remark}
\newtheorem{remarks}[subsection]{Remarks}

\numberwithin{equation}{subsection}

% Macros
%
\def\lim{\mathop{\mathrm{lim}}\nolimits}
\def\colim{\mathop{\mathrm{colim}}\nolimits}
\def\Spec{\mathop{\mathrm{Spec}}}
\def\Hom{\mathop{\mathrm{Hom}}\nolimits}
\def\Ext{\mathop{\mathrm{Ext}}\nolimits}
\def\SheafHom{\mathop{\mathcal{H}\!\mathit{om}}\nolimits}
\def\SheafExt{\mathop{\mathcal{E}\!\mathit{xt}}\nolimits}
\def\Sch{\mathit{Sch}}
\def\Mor{\mathop{\mathrm{Mor}}\nolimits}
\def\Ob{\mathop{\mathrm{Ob}}\nolimits}
\def\Sh{\mathop{\mathit{Sh}}\nolimits}
\def\NL{\mathop{N\!L}\nolimits}
\def\CH{\mathop{\mathrm{CH}}\nolimits}
\def\proetale{{pro\text{-}\acute{e}tale}}
\def\etale{{\acute{e}tale}}
\def\QCoh{\mathit{QCoh}}
\def\Ker{\mathop{\mathrm{Ker}}}
\def\Im{\mathop{\mathrm{Im}}}
\def\Coker{\mathop{\mathrm{Coker}}}
\def\Coim{\mathop{\mathrm{Coim}}}

% Boxtimes
%
\DeclareMathSymbol{\boxtimes}{\mathbin}{AMSa}{"02}

%
% Macros for moduli stacks/spaces
%
\def\QCohstack{\mathcal{QC}\!\mathit{oh}}
\def\Cohstack{\mathcal{C}\!\mathit{oh}}
\def\Spacesstack{\mathcal{S}\!\mathit{paces}}
\def\Quotfunctor{\mathrm{Quot}}
\def\Hilbfunctor{\mathrm{Hilb}}
\def\Curvesstack{\mathcal{C}\!\mathit{urves}}
\def\Polarizedstack{\mathcal{P}\!\mathit{olarized}}
\def\Complexesstack{\mathcal{C}\!\mathit{omplexes}}
% \Pic is the operator that assigns to X its picard group, usage \Pic(X)
% \Picardstack_{X/B} denotes the Picard stack of X over B
% \Picardfunctor_{X/B} denotes the Picard functor of X over B
\def\Pic{\mathop{\mathrm{Pic}}\nolimits}
\def\Picardstack{\mathcal{P}\!\mathit{ic}}
\def\Picardfunctor{\mathrm{Pic}}
\def\Deformationcategory{\mathcal{D}\!\mathit{ef}}
\newenvironment{reference}{}{}
\newenvironment{slogan}{}{}
\newenvironment{history}{}{}
\usepackage{cleveref}
\crefname{originallemma}{original lemma}{original lemmas}
\Crefname{originallemma}{Original lemma}{Original lemmas}
\makeatletter
\newlabel{lemma-flat-base-change-lci}{{069I}{source}{}{sourcepointer-1}{}}
\newlabel{lemma-flat-base-change-lci@cref}{{[originallemma][0][]069I}{[1][0][]source}{}{}{}}
\makeatother
\crefname{originallemma}{original lemma}{original lemmas}
\Crefname{originallemma}{Original lemma}{Original lemmas}
\makeatletter
\newlabel{more-morphisms-lemma-descending-property-lci}{{069N}{source}{}{sourcepointer-2}{}}
\newlabel{more-morphisms-lemma-descending-property-lci@cref}{{[originallemma][0][]069N}{[1][0][]source}{}{}{}}
\makeatother
\crefname{originallemma}{original lemma}{original lemmas}
\Crefname{originallemma}{Original lemma}{Original lemmas}
\makeatletter
\newlabel{more-morphisms-lemma-lci-syntomic-local-source}{{069P}{source}{}{sourcepointer-3}{}}
\newlabel{more-morphisms-lemma-lci-syntomic-local-source@cref}{{[originallemma][0][]069P}{[1][0][]source}{}{}{}}
\makeatother
\crefname{originallemma}{original lemma}{original lemmas}
\Crefname{originallemma}{Original lemma}{Original lemmas}
\makeatletter
\newlabel{morphisms-lemma-etale-syntomic}{{02GQ}{source}{}{sourcepointer-4}{}}
\newlabel{morphisms-lemma-etale-syntomic@cref}{{[originallemma][0][]02GQ}{[1][0][]source}{}{}{}}
\makeatother
\crefname{originallemma}{original lemma}{original lemmas}
\Crefname{originallemma}{Original lemma}{Original lemmas}
\makeatletter
\newlabel{spaces-topologies-definition-fpqc-covering}{{03MQ}{source}{}{sourcepointer-5}{}}
\newlabel{spaces-topologies-definition-fpqc-covering@cref}{{[originallemma][0][]03MQ}{[1][0][]source}{}{}{}}
\makeatother
\crefname{originallemma}{original lemma}{original lemmas}
\Crefname{originallemma}{Original lemma}{Original lemmas}
\makeatletter
\newlabel{spaces-topologies-lemma-recognize-fpqc-covering}{{03MS}{source}{}{sourcepointer-6}{}}
\newlabel{spaces-topologies-lemma-recognize-fpqc-covering@cref}{{[originallemma][0][]03MS}{[1][0][]source}{}{}{}}
\makeatother
\crefname{originallemma}{original lemma}{original lemmas}
\Crefname{originallemma}{Original lemma}{Original lemmas}
\makeatletter
\newlabel{spaces-topologies-lemma-refine-fpqc-schemes}{{0419}{source}{}{sourcepointer-7}{}}
\newlabel{spaces-topologies-lemma-refine-fpqc-schemes@cref}{{[originallemma][0][]0419}{[1][0][]source}{}{}{}}
\makeatother
\title{Being a local complete intersection morphism is fpqc local on the base for algebraic spaces}
\author{The Stacks Project Authors}
\date{stacks-project:spaces-more-morphisms.tex:local-snapshot-20261011; GNU FDL 1.2 or later}
\begin{document}
\maketitle
\noindent\textit{Editorial note.} Statement span is lines 11315-11319 of sources/stacks-spaces-more-morphisms/pinned-original-source.tex; proof span is lines 11321-11368, the complete author proof (46 body lines).\par\medskip
\noindent\textit{Editorial note.} Navigation stubs: 7 cited-but-not-included labels are supplied as non-load-bearing pointers with their pinned tags where the snapshot tag index has them.\par\medskip
\noindent\textit{Editorial note.} extra\_packages is empty on purpose: the retained preamble already loads xy with options and a second declaration clashes in the assembler wrapper.\par\medskip
\noindent\textit{Editorial note.} Text is retained verbatim, including the author's own typos and placeholders.\par\medskip
\noindent\textit{Editorial note.} Counted B side (both complete; spaces more general).\par\medskip
\noindent\textit{Allowed standard prerequisites (editorial).}\begin{itemize}
\item Explicit prerequisite boundary. More on Morphisms, lemma descending-property-lci (schemes side); More on Morphisms, lemma lci-syntomic-local-source (card L16201); Morphisms, lemma etale-syntomic; chapter lemma flat-base-change-lci; Topologies on Spaces chart lemmas.
\end{itemize}\medskip
\ifdefined\hypertarget\hypertarget{sourcepointer-1}{}\fi
\noindent\textit{Source reference.} Original lemma 069I: sources/stacks-spaces-more-morphisms/pinned-original-source.tex, source lines 11228--11231. Cited by the retained statement or proof but not part of the collected spans; supplied only as a navigation pointer. The argument does not depend on its content.\par\medskip
\ifdefined\hypertarget\hypertarget{sourcepointer-2}{}\fi
\noindent\textit{Source reference.} Original lemma 069N: sources/stacks-spaces-more-morphisms/pinned-original-source.tex, source lines 11368--11368. Cited by the retained statement or proof but not part of the collected spans; supplied only as a navigation pointer. The argument does not depend on its content.\par\medskip
\ifdefined\hypertarget\hypertarget{sourcepointer-3}{}\fi
\noindent\textit{Source reference.} Original lemma 069P: sources/stacks-spaces-more-morphisms/pinned-original-source.tex, source lines 11368--11368. Cited by the retained statement or proof but not part of the collected spans; supplied only as a navigation pointer. The argument does not depend on its content.\par\medskip
\ifdefined\hypertarget\hypertarget{sourcepointer-4}{}\fi
\noindent\textit{Source reference.} Original lemma 02GQ: sources/stacks-spaces-more-morphisms/pinned-original-source.tex, source lines 11368--11368. Cited by the retained statement or proof but not part of the collected spans; supplied only as a navigation pointer. The argument does not depend on its content.\par\medskip
\ifdefined\hypertarget\hypertarget{sourcepointer-5}{}\fi
\noindent\textit{Source reference.} Original lemma 03MQ: sources/stacks-spaces-more-morphisms/pinned-original-source.tex, source lines 11368--11368. Cited by the retained statement or proof but not part of the collected spans; supplied only as a navigation pointer. The argument does not depend on its content.\par\medskip
\ifdefined\hypertarget\hypertarget{sourcepointer-6}{}\fi
\noindent\textit{Source reference.} Original lemma 03MS: sources/stacks-spaces-more-morphisms/pinned-original-source.tex, source lines 11368--11368. Cited by the retained statement or proof but not part of the collected spans; supplied only as a navigation pointer. The argument does not depend on its content.\par\medskip
\ifdefined\hypertarget\hypertarget{sourcepointer-7}{}\fi
\noindent\textit{Source reference.} Original lemma 0419: sources/stacks-spaces-more-morphisms/pinned-original-source.tex, source lines 11368--11368. Cited by the retained statement or proof but not part of the collected spans; supplied only as a navigation pointer. The argument does not depend on its content.\par\medskip
\begin{lemma}
\label{lemma-descending-property-lci}
The property $\mathcal{P}(f) =$``$f$ is a local complete intersection
morphism'' is fpqc local on the base.
\end{lemma}

\begin{proof}
Let $S$ be a scheme. Let $f : X \to Y$ be a
morphism of algebraic spaces over $S$.
Let $\{Y_i \to Y\}$ be an fpqc covering
(Topologies on Spaces, Definition
\ref{spaces-topologies-definition-fpqc-covering}).
Let $f_i : X_i \to Y_i$ be the base change of $f$ by $Y_i \to Y$.
If $f$ is a local complete intersection morphism,
then each $f_i$ is a local complete intersection morphism
by Lemma \ref{lemma-flat-base-change-lci}.

\medskip\noindent
Conversely, assume each $f_i$ is a local complete intersection morphism.
We may replace the covering by a refinement (again because
flat base change preserves the property of being a
local complete intersection morphism). Hence we may assume
$Y_i$ is a scheme for each $i$, see
Topologies on Spaces, Lemma \ref{spaces-topologies-lemma-refine-fpqc-schemes}.
Choose a scheme $V$ and a surjective \'etale morphism $V \to Y$.
Choose a scheme $U$ and a surjective \'etale morphism
$U \to V \times_Y X$. We have to show that $U \to V$ is a
local complete intersection morphism of schemes.
By Topologies on Spaces, Lemma
\ref{spaces-topologies-lemma-recognize-fpqc-covering}
we have that $\{Y_i \times_Y V \to V\}$ is an fpqc covering
of schemes. By the case of schemes
(More on Morphisms, Lemma \ref{more-morphisms-lemma-descending-property-lci})
it suffices to prove the base change
$$
U \times_Y Y_i = U \times_V (V \times_Y Y_i) \longrightarrow V
$$
of $U \to V$ by $V \times_Y Y_i \to V$ is a
local complete intersection morphism. We can write this as the
composition
$$
U \times_Y Y_i \longrightarrow
(V \times_Y X) \times_Y Y_i = V \times_Y X_i \longrightarrow
V \times_Y Y_i
$$
The first arrow is an \'etale morphism of schemes (as a base change of
$U \to V \times_Y X$) and the second arrow is a
local complete intersection morphism of schemes as a flat base change of $f_i$.
The result follows as being a local complete intersection morphism
is syntomic local on the source and since \'etale morphisms
are syntomic (More on Morphisms, Lemma
\ref{more-morphisms-lemma-lci-syntomic-local-source}
and Morphisms, Lemma \ref{morphisms-lemma-etale-syntomic}).
\end{proof}

\end{document}
